\documentclass[11pt]{article}
\usepackage{ap-lang}
\title{Part 2: Claims and Evidence}
\author{Paul Wang}
\date{November 2026}
\docpart{Part 2 \textperiodcentered\ Study Guide}

\begin{document}
\maketitle

\begin{abstract}
An argument is a claim held up by evidence and tied to it by reasoning. This guide covers the kinds of claims and how a thesis differs from a topic (Section 1), the kinds of evidence and the jobs evidence does (Section 2), commentary, the part of an argument that explains the evidence and the part that earns most essay points (Section 3), qualification, concession, and rebuttal (Section 4), and how these ideas are tested in writing questions (Section 5). Section 6 is a method for the argument essay (Question 3). Each section ends with check questions; the key is on the last page.
\end{abstract}

\scoringline{Check questions: \blank\ / \qtotal \qquad Target: 13+ \qquad Terms I mix up: \blank[10em]}

\section{Claims and the thesis}

A \textbf{claim} is an assertion that can be argued: someone could reasonably disagree. The \textbf{thesis} is the central claim the whole text defends; \textbf{subclaims} (supporting claims) are the steps toward it.

\begin{compare}{Type}{Asserts\ldots}{Example}
Claim of fact & that something is or was so (and can be checked). & ``Teen sleep has declined since 2010.'' \\
Claim of value & that something is good, bad, right, or wrong. & ``Grading homework for completion rewards compliance over learning.'' \\
Claim of policy & that something should be done. & ``The district should start high school after 8:30 a.m.'' \\
Claim of cause & that one thing produces another. & ``Early start times, more than phones, explain teen sleep loss.'' \\
\end{compare}

\begin{keyidea}[Topic vs.\ thesis]
``Social media and teenagers'' is a topic. ``Social media and teenagers is a complex issue'' is still a topic, dressed up. ``Schools should teach students to curate their feeds rather than ban phones, because restriction ends at the school door and judgment does not'' is a thesis: it takes a position and signals a reason.
\end{keyidea}

\begin{q}{3}{2.1 Claims}{``Should'' + an action = claim of policy.}
``The city should convert two downtown parking lots into public housing'' is best described as a claim of
\chh{fact}{value}{policy}{cause}
\end{q}

\begin{q}{2}{2.1 Thesis}{A thesis takes an arguable position, not just a topic.}
Which sentence would work best as a thesis?
\ch{``Standardized testing has many pros and cons.''}{``Colleges should keep testing optional, because a single score rewards test preparation money more than readiness.''}{``This essay will discuss standardized testing.''}{``Standardized tests were first used in the early twentieth century.''}
\end{q}

\section{Evidence: types and functions}

\subsection{Types}
\begin{terms}
Data and statistics & Numbers from surveys, studies, records. Strong for claims of fact and cause; persuasive to skeptical or expert audiences. \\
Anecdote & A brief story, often personal. Makes an issue concrete and human; weak alone as proof of a general claim. \\
Example & A specific instance of a general claim. \\
Expert testimony & Quotation or findings from an authority. Borrows credibility. \\
Historical evidence & A precedent or past event. Shows a pattern or a consequence. \\
Analogy & A comparison to a parallel case. Explains or suggests; does not prove. \\
Observation & What the writer saw or experienced firsthand. \\
Hypothetical & An imagined scenario (``Imagine a town where\ldots''). Illustrates; weakest as proof. \\
\end{terms}

\subsection{Functions}
Reading questions usually ask what evidence \emph{does}. Learn the verbs:
\begin{terms}
Illustrate / exemplify & Give a concrete instance of an abstract claim. \\
Support / substantiate & Supply proof for the claim. \\
Amplify & Add weight or intensity to a point already made. \\
Contrast / juxtapose & Set two cases side by side to sharpen a difference. \\
Qualify & Limit the scope of a claim. \\
Refute & Disprove an opposing claim. \\
Establish credibility & Show the writer is informed or trustworthy. \\
\end{terms}

\begin{trap}
Anecdote as proof. An anecdote illustrates; it does not by itself establish a general claim. If a choice says a single story ``proves'' a trend, it is usually wrong.
\end{trap}

\begin{q}{4}{2.2 Evidence}{One story illustrates and humanizes; it cannot prove a general claim.}
A writer arguing that a highway expansion will divide a neighborhood tells the story of a grandmother who will no longer be able to walk to her church. The anecdote functions mainly to
\ch{prove that most residents oppose the expansion}{provide statistical evidence of traffic patterns}{refute the engineers' cost estimates}{make the abstract harm concrete and human}
\end{q}

\begin{q}{2}{2.2 Evidence}{An authority's findings borrow credibility for the claim.}
A writer quotes the findings of a state geologist to support the claim that a hillside is unstable. This evidence is best described as
\chh{a hypothetical}{expert testimony}{an analogy}{an anecdote}
\end{q}

\begin{q}{3}{2.2 Evidence}{An analogy explains by comparison but cannot prove.}
A writer compares the national debt to a household's credit card. The comparison is most effective as
\ch{a statistic measuring the debt}{proof that the debt will cause a recession}{a way of explaining an abstract idea in familiar terms}{a refutation of economists' views}
\end{q}

\section{Commentary}

\textbf{Commentary} is the writer's explanation of how evidence supports a claim: the \emph{because}, the \emph{which means}, the \emph{so what}. Evidence without commentary is a list; commentary without evidence is opinion.

\begin{study}{Commentary in action}
Evidence: ``After the district moved start times to 8:45, first-period tardiness fell by 30\%.''
No commentary: ``This shows tardiness fell.'' (restates.)
Commentary: ``Tardiness is the most visible symptom of exhaustion; when a schedule change alone erases nearly a third of it, the problem was never student laziness but a bell that rang before teenage bodies were awake.''
\end{study}

Good commentary does at least one of these: explains the \textbf{mechanism} (how the evidence leads to the claim), draws out an \textbf{implication}, rules out an \textbf{alternative explanation}, or connects back to the \textbf{thesis}.

\begin{q}{1}{2.3 Commentary}{Commentary explains how the evidence proves the claim; it does not restate it.}
Which sentence best provides commentary on the evidence ``Volunteer hours in the city fell 40\% after the community center closed''?
\ch{``The drop suggests the center did more than host programs: it was where neighbors found out who needed help.''}{``The community center was closed in 2021.''}{``Volunteering is important.''}{``Volunteer hours dropped by 40\%.''}
\end{q}

\begin{q}{2}{2.3 Commentary}{Ruling out an alternative explanation strengthens the causal claim.}
A writer notes that test scores rose after a school added recess, then adds: ``The scores rose only at the schools that added recess, not at the others in the district using the same new curriculum.'' This added sentence mainly
\ch{repeats the evidence}{rules out an alternative explanation for the rise}{introduces a new topic}{concedes that recess may lower scores}
\end{q}

\section{Qualification, concession, and rebuttal}

\begin{terms}
Qualifier & A word or phrase that limits a claim: \emph{often, most, in many cases, for now, at least in cities}. Makes a claim more defensible, not weaker. \\
Concession & Granting that part of the opposing view is true: ``Critics are right that the program is expensive.'' \\
Rebuttal (refutation) & Answering the opposing view: ``But the cost of doing nothing is higher.'' \\
Counterargument & The opposing view itself, stated so it can be answered. \\
\end{terms}

The common pattern is \textbf{concede and refute}: state the strongest version of the other side, grant what is true in it, then show why the claim still holds. It builds credibility (the writer is fair) and pre-empts the audience's objections.

\begin{trap}
A concession is not a change of position. When a writer says ``I will admit that\ldots,'' look for the ``but'' that follows; the concession usually sets up a stronger version of the original claim.
\end{trap}

\begin{q}{3}{2.4 Qualification}{Limiting the claim's scope makes it more defensible.}
A writer revises ``Homework is useless'' to ``Homework beyond an hour a night adds little for most middle school students.'' The revision
\ch{introduces a counterargument}{abandons the writer's position}{qualifies the claim, making it more defensible}{adds an anecdote}
\end{q}

\begin{q}{4}{2.4 Concession}{Granting a point then answering it = concession and rebuttal.}
``Opponents are right that the stadium will bring jobs. But most of those jobs last only as long as construction, while the tax break lasts thirty years.'' The passage
\ch{uses an analogy to explain the stadium's cost}{rejects every claim made by opponents}{changes the writer's position on the stadium}{concedes a point and then rebuts it}
\end{q}

\section{Claims and evidence in writing questions}

Writing questions test the same skills from the writer's side:
\begin{itemize}[leftmargin=1.5em]
\item \textbf{Best evidence}: the choice that supports \emph{the specific claim in that sentence}, not the topic in general. Check scope: a national statistic does not support a claim about one school.
\item \textbf{Qualify a claim}: the choice that limits an overstatement without abandoning it.
\item \textbf{Counterargument}: the choice that states an objection the audience might actually hold, then sets up the answer.
\item \textbf{Relevance}: ``Should the writer keep this sentence?'' Keep it only if it advances the paragraph's claim.
\end{itemize}

\begin{q}{4}{2.5 Writing: evidence}{Evidence must match the specific claim (here, this school's lunch period).}
A draft claims: ``At Lincoln High, the 25-minute lunch period is too short for students to eat.'' Which evidence best supports this claim?
\ch{Some schools in Europe have hour-long lunches.}{National guidelines recommend that students have time to eat.}{Many students enjoy pizza for lunch.}{A survey of Lincoln students found that 61\% regularly leave lunch with food uneaten because the line takes 15 minutes.}
\end{q}

\section{Writing the argument essay (Q3)}

\subsection{The prompt}
Q3 gives a short quotation or idea and asks you to ``write an essay that argues your position on'' it. There are no sources. Your evidence comes from what you know: history, current events, science, literature, the arts, your community, and your own experience.

\subsection{Choosing a position}
Choose the position you can support with the most \emph{specific} evidence, not the one you find most admirable. A qualified position (``yes, but only when\ldots'') is often the most defensible and leads naturally to Row C.

\subsection{Evidence that earns Row B}
\begin{enumerate}[leftmargin=1.5em]
\item \textbf{Specific over general}: ``the 1965 Selma marches'' beats ``protests in history.''
\item \textbf{Two or three developed examples} beat six mentioned ones.
\item \textbf{Personal experience is allowed} and can be strong if specific and analyzed, not just told.
\item \textbf{Every example gets commentary} that ties it back to the claim of its paragraph.
\end{enumerate}

\subsection{Reaching Row C}
Explore the \textbf{complexity}: when does your claim stop being true? What tension does the prompt contain (freedom vs.\ security, individual vs.\ community)? Situate the argument in a broader context, or show why the question matters now.

\begin{q}{1}{2.6 Argument essay}{Specific, developed evidence with commentary moves Row B.}
Which plan is most likely to earn a high Row B score on an argument essay about whether failure is necessary for success?
\ch{Two developed examples, a scientist's failed experiments and the writer's own failed tryout, each followed by commentary explaining what the failure taught}{Six one-sentence examples of famous people who failed}{A hypothetical scenario about an imaginary inventor}{A summary of the prompt followed by the writer's opinion}
\end{q}

\begin{q}{1}{2.6 Argument essay}{A qualified position explores complexity, the path to Row C.}
Which thesis best positions an argument essay to earn the sophistication point?
\ch{``Failure helps people succeed when it comes with the time and support to learn from it, which is why who can afford to fail matters as much as failure itself.''}{``Failure is bad and should be avoided.''}{``Everyone fails sometimes.''}{``There are many opinions about failure.''}
\end{q}

\makekey
\end{document}
